\section{Introduction}

% [1] https://science.nasa.gov/dark-energy/#:~:text=Astronomer%20Edwin%20Hubble%20confirmed%20that,they%20confirmed%2C%20is%20really%20expanding.
% [1] https://www.pnas.org/doi/10.1073/pnas.15.3.168
% [2] https://arxiv.org/abs/astro-ph/9712212
% [3] https://arxiv.org/abs/2203.06142v1

% [1] https://arxiv.org/abs/astro-ph/9805201
The expansion of the Universe was confirmed in 1929 through observations of the redshifts of spiral galaxies (\cite{hubbleRelationDistanceRadial1929a}). At the time, it was expected that gravitational attraction would act to decelerate this expansion, a prediction consistent with the framework of general relativity for a matter-dominated Universe. However, observations of high-redshift Type Ia supernovae revealed them to be further and fainter than expected, providing the first direct evidence that the expansion of the Universe is in fact accelerating (\cite{perlmutterDiscoverySupernovaExplosion1998a}, \cite{riessObservationalEvidenceSupernovae1998}). This discovery significantly changed our understanding of the Universe. In response, the scientific community turned to a range of independent observational probes, such as the Cosmic Microwave Background (CMB), Type Ia Supernovae, and Baryon Acoustic Oscillations, to better understand this acceleration (\cite{abdallaCosmologyIntertwinedReview2022a}). 
% also talk about any information we can get from measuring the acceleration of the Universe and why it's significant
% dark matter and dark energy
% talk about the current standard cosmological model

% [4] https://arxiv.org/html/2312.08542v1?st_source=ai_mode#:~:text=One%20of%20the%20most%20important,et%20al.%2C%202022)%20.
% [5] https://arxiv.org/abs/1807.06209
% [6] https://arxiv.org/abs/2112.04510
% [7] https://arxiv.org/abs/2106.15656
% [8] https://arxiv.org/abs/2308.10954
% [9] https://arxiv.org/abs/2107.10291
% [10] https://arxiv.org/abs/2103.01183
Since this breakthrough, as instrumental technology advanced, the volume and precision of cosmological data increased significantly and led to high precision in cosmological parameters. While the majority of observations agree with the predictions of the standard cosmological model, there are discrepancies between these methods, known as cosmological tensions (\cite{leizerovichTensionsCosmologyDiscussion2024a}). The most prominent being the discrepancy between the Hubble constant $H_0$ determined from the CMB (\cite{collaborationPlanck2018Results2020a}) not being in agreement with $H_0$ determined from Type Ia supernovae and Cepheids (\cite{riessComprehensiveMeasurementLocal2022a}). Although this discrepancy has been discussed at length in literature, there is no agreement in the scientific community on the source (\cite{freedmanMeasurementsHubbleConstant2021}, \cite{riessLocalValueH$_0$2024}, \cite{schoneberg$H_0$OlympicsFair2022}). Due to the increased precision in the current cosmological probes becoming close to reaching its maximum, it is prudent to look for new and independent cosmological probes that could either confirm or deny the discrepancies found (\cite{valentinoRealmHubbleTension2021}).

% [11] https://www.cfa.harvard.edu/research/topic/large-scale-structure?st_source=ai_mode (first sentence but havent cited it yet)
% [12] https://arxiv.org/pdf/2206.01579
% or https://academic.oup.com/mnras/article/423/4/3430/1747727?st_source=ai_mode ??
% [13] https://academic.oup.com/mnras/article/355/4/1339/993455?login=false
% [14] https://academic.oup.com/mnras/article/541/1/476/8161700
Probing the large scale structure (LSS) of the Universe is crucial in providing information about the very early Universe, the distribution of dark matter, galaxy formation, and measuring the Universe's acceleration. Traditionally, this involved measuring the redshift and angular position of galaxies individually, then creating a catalogue and 3D map in the optical or near-infrared (\cite{cunningtonHIIntensityMapping2022}). While this method led to constraints on dark energy, gravity, and the initial conditions of the Universe, it often requires the difficult detection of individual galaxies which becomes expensive. An alternative and novel probe for the LSS of the Universe is Neutral Hydrogen (HI) Intensity Mapping (\cite{battyeNeutralHydrogenSurveys2004}) which overcomes these challenges using the cumulative emission of HI flux in the sky (\cite{mazumderHiIntensityMapping2025}). Hence, it can map large areas of the sky very quickly. 

% [14] https://arxiv.org/pdf/0902.3091
% [15] https://arxiv.org/pdf/0801.1677
% [16] https://iopscience.iop.org/article/10.1086/504972
% [17] https://journals.aps.org/prd/abstract/10.1103/PhysRevD.76.083005
Following the epoch of reionization, majority of HI exists inside of galaxies due to self-shielding. This makes HI an excellent tool for probing how matter is distributed throughout the Universe. The 21\,cm line arising from the hyperfine spin-flip transition of HI is the dominant spectral line in radio astronomical observations at frequencies below 1420MHz (\cite{peterson21CmIntensity2009}). By observing this redshifted line with radio telescopes, and due to the fact it is an isolated transition, a 3D map of the LSS of the Universe can be constructed providing the radial coordinate along the line-of-sight through the frequency dependence of the redshift. This is a significant advantage over probes such as the CMB, which maps a single 2D surface (\cite{loebPossibilityPreciseMeasurement2008a}). Furthermore, this line can in principle be observed across a wide range of cosmic epochs, from the dark ages through the epoch of reionisation to the present day (\cite{shapiro21CmBackground2006}), \cite{lewis21CmAngularpower2007a}), enabling access to the large co-moving volumes at high redshift where the majority of the Universe's structure resides. HI intensity mapping therefore offers a powerful probe of the large-scale matter distribution.

% the faintness of HI emission lines make it challenging to resolve individual galaxies but for cosmology we are interested in the 
% talk about meerkat arrays. mapping out stuff with sufficient baseline/sensitivity

% [18] https://ui.adsabs.harvard.edu/abs/1999A%26A...345..380S/abstract
% [19] https://arxiv.org/pdf/astro-ph/0302099 
One of the main challenges for HI intensity mapping is the presence of foreground emission, such as synchrotron and free-free emission from the Milky Way (\cite{shaverCanReionizationEpoch1999a}), which dominates at this wavelength. The foregrounds are expected to be spectrally and spatially smooth in comparison to the rapid variations in frequency of the HI signal, which means that even though the HI signal is several orders of magnitude fainter than the foreground, the signal can be separated out (\cite{ohForegrounds21cmObservations2003a}). 

% [2] https://academic.oup.com/mnras/article/488/4/5452/5531337
% [3] https://research.manchester.ac.uk/en/publications/interference-from-global-navigation-satellites-in-future-hi-inten/
Beyond foreground contamination, HI intensity mapping faces additional challenges including thermal noise from the receiver system, frequency-dependent beam effects that cause mixing between angular and spectral modes (\cite{cunningtonImpactForegroundsIntensity2019}), and radio frequency interference from human-made transmitters (\cite{harperInterferenceGlobalNavigation2018}). Current and next-generation radio telescopes such as MeerKat and the Square Kilometre Array (SKA; \cite{groupCosmologyPhase12020}) are expected to map HI intensity across large volumes of the Universe with unprecedented sensitivity (\cite{hothiComparingForegroundRemoval2021}), making the development and validation of robust foreground removal techniques an increasingly urgent priority.
% [7] https://arxiv.org/abs/1811.02743
% [8] https://arxiv.org/abs/1501.03989

% [4] https://arxiv.org/abs/2511.13328?st_source=ai_mode#:~:text=Our%20goal%20is%20to%20evaluate,%5Bastro%2Dph.CO%5D
% [5] https://iopscience.iop.org/article/10.1086/310631
% [6] https://arxiv.org/abs/2010.15843
% [8] https://academic.oup.com/mnras/article/500/2/2264/5979792
To address these challenges, a range of foreground removal techniques have been proposed and tested in the literature, including polynomial fitting along the frequency axis (\cite{gkogkouForegroundRemovalHI2026}), Independent Component Analysis (\cite{hothiComparingForegroundRemoval2021}), Wiener filtering (\cite{tegmarkHowMakeMaps1997}), and more recently deep learning approaches (\cite{makinenDeep21DeepLearning2021}). Each method represents a trade-off between the assumptions required and the accuracy of the recovered signal.

This work applies and compares two foreground removal methods: Principal Component Analysis (PCA) and Gaussian Process Regression (GPR). PCA is a blind statistical method which exploits the spectral smoothness of the foregrounds in comparison to the rapid fluctuations of the HI signal. GPR is a method that explicitly models the signal covariance. These methods are applied to simulated 21\,cm intensity mapping data with performance being evaluated through the recovery of the HI power spectrum, in order to assess the relative strengths and limitations of each approach.

The paper is organised as follows: \Cref{sec:simulated_dataset} introduces the simulated dataset. \Cref{sec:foreground_removal_methods} describes the PCA and GPR foreground removal methodologies. \Cref{sec:results} presents the recovered HI power spectrum results for both methods. \Cref{sec:conclusion} provides a summary and conclusions.
